\section{Introduction}
\label{sec:intro}

Mixture-of-Experts (MoE) increases model capacity through conditional computation, but this advantage is realized only when routing induces meaningful expert specialization \citep{fedus2022switch}. If routed experts remain functionally redundant, adding experts increases architectural complexity without converting that complexity into effective capacity. The central question is therefore not sparse activation by itself, but how gating produces -- or fails to produce -- specialized experts.

Our current study starts from the same broad concern as recent spectral analyses of MoE specialization failure: understanding why experts can remain similar even when routing appears non-trivial. However, our present checkpoint leads to a different diagnosis. In our setting, the most stable phenomenon is not a clean dominant-parameter-collapse pattern. Instead, we consistently observe a strong shared/common input subspace together with only very limited downstream gains from our current routing variants. \textcolor{red}{In particular, the current centered, ACP-style, and quadratic routers do not yet improve enough over a standard linear router with load balancing to support a final method claim in the introduction.}

This shifts the mainline of the paper. Our current evidence suggests that increasing geometric separation in representation space is not sufficient for expert specialization. A router can look more geometry-aware, or align better with certain expert subspaces, while the experts still fail to receive genuinely differentiated updates. In other words, \emph{representation-level separation does not automatically become update-level specialization}. We regard this gap between routing geometry and expert learning dynamics as the root problem exposed by the current milestone.

The anchor of our analysis remains the mean bias and anisotropy of language-model representations \citep{ethayarajh2019contextual,chi2022representation}. Hidden states contain a strong common component, and routing without mean handling can easily be dominated by that shared direction. Centering therefore helps expose residual variation that is more relevant for expert discrimination. But our current results also indicate that centering alone is not the whole story: even when the geometry of the routed representation becomes more separated, the experts themselves do not necessarily become more specialized under spectral audits. The unresolved question is how the representation-side common mode, the gate, and the resulting expert gradients are coupled.

Accordingly, the goal of this work is no longer to prematurely commit to one routing form, but to clarify what kind of gating can actually induce expert specialization. We study this question from the perspective of spectral geometry: how shared directions in the input representation affect routing, how routing affects the distribution of expert updates, and what measurements are needed to distinguish superficial geometric separation from useful specialization. \textcolor{red}{The final routing mechanism of the paper is therefore still open at the current stage; the role of the present draft is to establish the problem, the observations, and the criteria that a successful design must satisfy.}

At the current stage, the paper is anchored by the following points:
\begin{itemize}
    \item \textbf{Expert specialization should be defined through the relation between routing and expert updates, rather than through occupancy or geometric dispersion alone.} Our current checkpoint suggests that better-looking separation in representation space does not automatically imply specialized experts.
    \item \textbf{Mean bias and anisotropy provide a principled anchor for studying routing failure.} Centering is important because a strong common component can dominate gating, but centering alone does not yet resolve the specialization problem.
    \item \textbf{The next step is to design gating by targeting update differentiation.} \textcolor{red}{Our current routing variants provide useful evidence and negative results, but they have not yet reached the root cause of expert specialization failure.}
\end{itemize}
